\section{Introduction}
\label{sec:introduction}

The rapid proliferation of shared electric scooters (e-scooters) has introduced a new class of urban safety challenge. Emergency department visits related to e-scooter injuries have surged globally, with head trauma and fractures dominating presentations \citep{trivedi2019injuries, bloom2021standing, badeau2019emergency}. Fatalities, though less frequent, often involve collisions with motor vehicles at speed \citep{yang2020safety}. Speed is widely recognized as the primary modifiable risk factor: the power model of the speed--crash relationship predicts that even modest speed increases lead to disproportionate injury risk, a relationship especially consequential for unprotected riders \citep{elvik2013reexamination}.

Recognizing this, jurisdictions worldwide have imposed regulatory speed limits on e-scooters---25~km/h in South Korea and France, 20~km/h in Germany and the UK, and 15--25~km/h across US cities \citep{ma2021municipal, sexton2023shared}. Yet a critical gap persists between the existence of speed limits and their enforcement. Unlike motor vehicles, which have physical governors linked to fuel systems, e-scooter compliance depends entirely on software-based speed governors embedded in the vehicle firmware. Operators can configure these governors to permit different maximum speeds through selectable riding modes. Whether these software governors actually reduce speeding in practice---and whether their removal leads to predictable increases in unsafe behavior---remains an open empirical question with direct policy implications.

Prior empirical work on e-scooter speed behavior has been constrained in three important respects. First, existing studies are small in scale: epidemiological studies document injury patterns but lack behavioral data on the rides that produced those injuries \citep{trivedi2019injuries, bloom2021standing}; naturalistic riding studies provide within-trip speed profiles but are limited to small samples (typically $N < 500$ trips from a single city) \citep{ma2021escooter}; and roadside observation studies capture spot speeds but not continuous profiles \citep{cicchino2023speed}. Second, no study has provided \textit{causal} evidence on speed governor effectiveness. Cross-sectional comparisons between governed and ungoverned modes are suggestive but cannot rule out self-selection: riders who choose ungoverned modes may differ systematically from those who choose governed modes. Third, behavioral compensation (risk homeostasis) after governor introduction or removal has not been tested---it is unknown whether riders compensate for the loss of an ungoverned option by speeding more aggressively on governed modes.

We address all three gaps using a uniquely comprehensive dataset. Swing, one of South Korea's largest shared e-scooter operators, offered three speed modes on identical hardware during 2023: TUB (turbo, no speed governor), STD (standard, governed to $\sim$25~km/h), and ECO (economy, governed to $\sim$20~km/h). Using 21~million GPS-instrumented trips with per-point sensor speed data from 52 cities over 10 months (February--November 2023), we conduct three independent analyses that together form a converging evidence funnel:

\begin{enumerate}
    \item \textbf{Cross-sectional evidence.} A logistic regression with generalized estimating equations (GEE) on the full dataset reveals that TUB mode increases speeding odds 121-fold relative to STD ($p < 0.001$), dominating all other predictors. A LightGBM model with SHAP interpretability (AUC $= 0.905$) confirms that operating mode is 3.6 times more important than any other feature. An experience analysis reveals that the positive experience--speeding gradient is almost entirely mediated (98.4\%) by progressive TUB adoption rather than within-mode behavioral changes.

    \item \textbf{Causal evidence.} In December 2023, Swing banned TUB mode system-wide under regulatory pressure. We exploit this exogenous shock in a difference-in-differences (DiD) design, using cross-city variation in pre-ban TUB adoption as continuous treatment intensity. Cities with higher pre-ban TUB share experienced proportionally larger speeding reductions ($\beta = -0.568$, $p < 0.001$; dose-response $r = 0.754$), with a null placebo test ($p = 0.56$) supporting causal validity. Critically, multi-outcome DiD and within-user mode-switcher analysis confirm no behavioral compensation: STD/ECO speeding rates are unchanged after the ban ($p = 0.59$; Cohen's $d = 0.069$).

    \item \textbf{Survival evidence.} Among 864,000 users with at least two trips, Kaplan-Meier analysis shows that TUB users reach their first speeding event after a median of 15 trips versus 81 trips for non-TUB users (5.4-fold). A Cox proportional hazards model with time-varying covariates estimates a hazard ratio of 10.41 for TUB adoption ($p < 0.001$).
\end{enumerate}

These converging findings carry three direct policy implications. First, \textit{mandatory speed governors} should be required on all shared e-scooters---the 121-fold difference in speeding odds is far larger than any other intervention lever. Second, \textit{graduated mode access} policies should restrict new riders to governed modes, as the experience--speeding pathway operates primarily through TUB adoption. Third, \textit{context-sensitive geofencing} can further reduce risk by imposing stricter speed limits near schools, hospitals, and pedestrian zones. To our knowledge, this is the largest GPS-based e-scooter speed study to date and the first to provide both causal evidence and behavioral compensation tests for speed governor effectiveness.
